% Generated draft: requires mathematical and source review.
\documentclass[11pt]{article}
\usepackage[margin=25mm]{geometry}
\usepackage{fontspec}
\setmainfont{Latin Modern Roman}
\usepackage{amsmath,amssymb,mathtools,unicode-math}
\setmathfont{Latin Modern Math}
\usepackage{enumitem,needspace,xurl,hyperref}
\setlength{\parindent}{0pt}
\setlength{\parskip}{5pt}
\setlength{\emergencystretch}{3em}
\newcommand{\N}{\mathbb{N}}
\newcommand{\Z}{\mathbb{Z}}
\newcommand{\Q}{\mathbb{Q}}
\newcommand{\R}{\mathbb{R}}
\newcommand{\C}{\mathbb{C}}
\newcommand{\F}{\mathbb{F}}
\newcommand{\A}{\mathbb{A}}
\newcommand{\PP}{\mathbb{P}}
\newcommand{\E}{\mathbb{E}}
\newcommand{\eps}{\varepsilon}
\newcommand{\dd}{\,\mathrm{d}}
\begin{document}
\section*{UnsolvedMath ID: 30001097}
\textbf{Source problem code:} catalogue entry
\subsection{Definitions and mathematical statement}
The source concerns the averaged Whitham equation for the action vector $I$ arising in the small-noise limit of a randomly perturbed KdV equation. It asks whether this averaged equation has a unique solution and a unique stationary measure.

In conventional action notation, $I=(I_j)_{j\geq1}$ is a nonnegative sequence in an appropriate phase space $\mathcal I$. A schematic stochastic averaged equation would be
\[
dI=\overline F(I)\,d\tau+\overline\sigma(I)\,dW_\tau,
\]
but the actual coefficients, noise realization, sequence space and time scaling are not provided. This display is not a reconstructed equation from the report.

The first question must be read for fixed initial data or initial law: for every admissible such datum, is there a solution in the intended class, and is it unique in the intended sense? Pathwise uniqueness for a fixed driving noise, uniqueness in law and uniqueness of a martingale problem are different assertions, and the record does not select one. It does not ask for all solutions with different initial data to be equal.

If the solution theory defines a Markov semigroup $P_t$ on $\mathcal I$, a stationary probability measure is a probability measure $\pi$ with
\[
\pi P_t=\pi\quad\text{for every }t\geq0.
\]
The second question asks whether such a measure exists and is unique in the intended admissible class. Without a specified unique solution law, stationarity would instead need definition for the relevant solution or martingale problem. Uniqueness of an initial-value solution does not by itself supply uniqueness of stationary measures, so both source tasks are retained separately.

The perturbed KdV equation, spatial domain and boundary conditions, forcing and damping, action space, averaged drift and diffusion, solution class and moment or support conditions on stationary measures are all absent. No finite-dimensional truncation, standard Brownian coefficient or fixed-noise theorem is substituted for the requested limiting problem.

The supplied triage checked 2026-08-21 labels the problem open. It reports subsequential Whitham limits and uniqueness in restricted cases from averaging results, while calling general uniqueness of solutions and stationary measures unresolved. This is dated source curation. The corrected-verified note reports source checking and normalization of the question.
\subsection{Short English statement}
For the stochastic averaged action equation from perturbed KdV, ask for well-posedness at fixed initial data and for exactly one stationary measure. The equation, state space and senses of uniqueness are not supplied and must be distinguished.
\subsection{Sources}
\begin{itemize}
\item Local UnsolvedMath registry, permanent ID 30001097, source code catalogue entry; supplied complete record and corrected-verified source-check note. Literature checked 2026-08-21: the assessment is supplied dated curation, not independently verified here.
\item Infinite Dimensional Random Dynamical Systems and Their Applications (2009), Oberwolfach Reports, DOI 10.4171/owr/2008/50, pp. 2815--2874: \url{https://doi.org/10.4171/owr/2008/50}. Proposed-year metadata gives 2008.
\item Kuksin--Piatnitski averaging: \url{https://arxiv.org/abs/0710.3869}.
\item Kuksin effective equation: \url{https://arxiv.org/abs/1002.1294}.
\end{itemize}
\end{document}
